% ===========================================================================
\section{Introduction}
\label{sec:introduction}
% ===========================================================================

Additive manufacturing (AM) supply chains increasingly rely on
permissioned blockchains for tamper-evident lifecycle records spanning
material certification, printing, inspection, and final
certification~\cite{alkhader2020blockchain, kurpjuweit2021blockchain}.
Prior work, including our unified provenance--reputation
chaincode~\cite{anonymous2026unified}, shows such records can be anchored
cheaply and coupled atomically to on-chain trust scoring.  These systems
share a rarely examined assumption: that the data arriving at the
chaincode interface faithfully describes what happened on the factory
floor.

This assumption is an instance of the \emph{oracle
problem}~\cite{caldarelli2020oracle}: a blockchain guarantees the integrity
of data \emph{after} commitment, but says nothing about the process that
produced the data.  In AM the stakes are concrete: sabotage attacks that
alter process parameters or design files produce parts that pass visual
inspection yet fail in service~\cite{belikovetsky2017dr0wned,
yampolskiy2018security}, and a provenance ledger fed by an unauthenticated
HTTP POST records such an attack as a valid lifecycle event.  The
trustworthiness of the on-chain record therefore reduces to that of the
\emph{physical-to-chain bridge}---the software and hardware path from the
printer's sensors to the ledger commitment.

This paper studies that bridge directly, with a working implementation on
real hardware rather than simulation.  We make four contributions:

\begin{enumerate}
  \item \textbf{A threat-model-driven trust-tier analysis
  (Section~\ref{sec:threat_model})} organizing bridge designs into four
  tiers---T0 (na\"ive submission), T1 (edge-signed telemetry bundles), T2
  (on-chain device registry), T3 (hardware root of trust)---stating
  explicitly which attacks each tier defeats and which it does not.  T1
  makes evidence tamper-evident after signing but cannot detect
  fabrication at source by an adversary holding the device key; we treat
  this boundary as a first-class result.

  \item \textbf{A canonical, device-signed telemetry bundle format
  (Section~\ref{sec:architecture})} binding, in one signed unit, the
  certified design (G-code SHA-256), the commanded process envelope, the
  observed telemetry, a derived adherence assessment, and the job
  outcome.  Canonical JSON makes the bundle hash reproducible by any
  verifier; that hash is the only artifact stored on-chain.

  \item \textbf{An end-to-end T1 implementation on real hardware
  (Section~\ref{sec:implementation}):} a Creality Ender-3, OctoPrint on a
  Raspberry Pi, an Ed25519 device key, and Hyperledger Fabric v3.1.0
  running our three-contract unified chaincode, evaluated on telemetry
  captured from \emph{five real print jobs}---three geometries and three
  outcome types, including a controlled thermal-anomaly injection---not
  synthetic data.

  \item \textbf{Physical adherence as a reputation signal
  (Section~\ref{sec:arch_integration}):} the assessment maps thermal
  adherence, design-hash compliance, and timeliness onto the $[0,1]$
  rating consumed by an on-chain Bayesian reputation engine, so one
  atomic transaction records the provenance event \emph{and} updates the
  operator's per-dimension trust posterior.
\end{enumerate}
